\section{Introduction}

In this paper, we show that repeatedly deploying large language models (LLMs)
and fine-tuning them on curated data from earlier deployments significantly
improves their planning capabilities.  This curation can be simply done by
\emph{validating} traces from previous generations, and \emph{selecting}
appropriate valid traces for future training. This mechanism produces a training
process conceptually similar to RL fine-tuning but where the reward signal is
left implicit.
The core idea is simple:
repeated deployment starts by users generating text with an LLM after its release. These texts go through a curation process, e.g.\ texts that do not capture user intent are rejected. Remaining texts are then shared to the web, with scrapes of the web including the curated text used to fine-tune the next generation of the LLM.

Iterative deployment is not a contrived setting: GPT-3.5 was trained on data
scraped from the web following GPT-3's deployment, with data shared on the web
at the time including curated texts generated by users using GPT-3
\cite{brown-et-al-neurips2020}. Similarly, GPT-4 was trained on data shared by
users from GPT-3.5 and GPT-3 \cite{openai-arxiv2024}, and so on. With agent
workflows becoming more commonplace, future training data will include agent
traces from prior model generations, similarly leading to an iterative training
on previous generation data. Figure~\ref{fig:pipeline} illustrates the
basics of this mechanism.



\begin{figure}[t]
\centering
\newcommand{\iconsize}{0.1\linewidth} %

\tikzset{
  fig/.style = {inner sep=0pt, align=center},
  caption/.style = {font=\small, text width=\iconsize, align=center},
  ->/.style = {-{Latex[length=3mm,width=2mm]}, line width=0.6pt},
}

\begin{tikzpicture}[node distance=1.9cm and 1.6cm]

\node[fig] (dataset) {
  \includegraphics[width=\iconsize]{figures/dataset-icon.png}\\
  \parbox{0.18\linewidth}{\centering \textbf{Input Prompts}}
};

\node[fig, right=1.2cm of dataset] (llm1) {
  \includegraphics[width=\iconsize]{figures/llm-icon.png}\\
  \parbox{0.18\linewidth}{\centering \textbf{LLM Gen. $n$}}
};

\node[fig, right=1.2cm of llm1] (traces) {
  \includegraphics[width=\iconsize]{figures/traces-icon.png}\\
  \parbox{0.18\linewidth}{\centering \textbf{Plans + Traces}}
};

\node[fig, right=1.2cm of traces] (val) {
  \includegraphics[width=\iconsize]{figures/val-icon.png}\\
  \parbox{\iconsize}{\centering \textbf{Validator}}
};


\node[fig, below=2.5cm of val] (correct) {
  \includegraphics[width=\iconsize]{figures/correct-traces-icon.png}\\
  \parbox{0.25\linewidth}{\centering \textbf{Valid Plans + Traces}}
};

\node[fig, left=2.5cm of correct] (finetune) {
  \includegraphics[width=\iconsize]{figures/fine-tuning-icon.png}\\
  \parbox{0.18\linewidth}{\centering \textbf{Supervised Fine-Tuning}}
};

\node[fig, below=0.2cm of traces] (prev-gen) {
  \includegraphics[width=\iconsize]{figures/correct-traces-icon.png}\\
  \parbox{0.25\linewidth}{\centering \textbf{Valid Plans + Traces From Previous Generations}}
};

\node[fig, below=2.5cm of dataset] (llm2) {
  \includegraphics[width=\iconsize]{figures/llm-icon.png}\\
  \parbox{0.18\linewidth}{\centering \textbf{LLM Gen. $n+1$}}
};

\draw[->] (dataset.east) -- (llm1.west);
\draw[->] (llm1.east) -- (traces.west);
\draw[->] (traces.east) -- (val.west);


\draw[->, shorten >=0.1pt, shorten <=0.1pt]
  ($ (val.south) + (0,-1mm) $)
    to[out=-90, in=90, looseness=1.0]
  ($ (correct.north) + (0,3mm) $);
\draw[->] (prev-gen.west) --  (finetune.north);
\draw[->] (correct.west) -- (finetune.east);
\draw[->] (finetune.west) -- (llm2.east);

\end{tikzpicture}
\caption{Single iteration of the iterative deployment mechanism for
  planning. Using a fixed a set of planning tasks, we prompt the current version
  of the LLM --- referred to as the generation $n$ of the model --- to solve
  these tasks. An external validator (e.g. a human using a chatbot, or a
  computer programme in the case of planning) identifies the tasks solved
  correctly. Their traces and plans, together with the traces and plans from tasks solved by
  previous generations, are then used to fine-tune generation $n$, producing
  generation $n+1$ of the model.}
\label{fig:pipeline}
\end{figure}


\pgfplotsset{compat=1.11,
    /pgfplots/ybar legend/.style={
    /pgfplots/legend image code/.code={%
       \draw[##1,/tikz/.cd,yshift=-0.25em]
        (0cm,0cm) rectangle (10pt,0.8em);},
   },
}

\begin{figure}[t]
\centering
\begin{tikzpicture}
  \begin{axis}[
   ybar=0pt,
   bar width=20pt,
   width=13cm,
   height=4cm,
   enlarge x limits=0.35,
   enlarge y limits={upper,value=0.08},
   ymin=0,
   ymax=220,
   ylabel={\# solved tasks},
   symbolic x coords={Blocksworld,Rovers,Sokoban},
   xtick=data,
   xtick align=inside,
   x tick label style={font=\small},
   xtick style={draw=none},
   ytick style={draw=none},
   axis x line*=bottom,
   axis y line*=left,
   major y tick style = {black},
   grid=major,
   grid style={dashed,gray!30},
   legend style={at={(0.5,-0.18)}, anchor=north,legend columns=4, /tikz/every even column/.append style={column sep=6pt}},
   ylabel style={yshift=-2mm},
   cycle list name=exotic,
   nodes near coords,
   nodes near coords style={font=\footnotesize, /pgf/number format/fixed},
   every node near coord/.append style={yshift=2pt},
    ]

  \addplot[
    pattern=north east lines,
    pattern color=darkNavy,
    draw=black,
  ] coordinates {
    (Blocksworld,52)
    (Rovers,41)
    (Sokoban,32)
  };
  \addplot[draw=black, fill=specialRed] coordinates {
    (Blocksworld,109)
    (Rovers,131)
    (Sokoban,58)
  };
  \addplot[draw=black, fill=specialOrange] coordinates {
    (Blocksworld,132)
    (Rovers,201)
    (Sokoban,72)
  };

  \addplot[draw=black, fill=specialYellow] coordinates {
    (Blocksworld,154)
    (Rovers,205)
    (Sokoban,96)
  };

\legend{
  Base,
  Gen. 1,
  Gen. 2,
  Gen. 5
}

  \end{axis}
\end{tikzpicture}
\caption{Summary of our main results. Number of solved tasks (with 1000 tasks
  per domain) for three different domains when comparing the base model with
  later deployed generations (generations 1, 2, and 5). Average over three
  separate runs. In all domains, the fifth generation more than doubles the
  performance of the base model.}
\label{fig:results-summary}
\end{figure}



We evaluate this mechanism in the well controlled setting of classical planning.
Iterative deployment captures a pattern common to \emph{planning}: when users use LLMs e.g. to review a product, help solve a reasoning task, or plan a trip, they are more inclined to share LLM results publicly if they are correct. This works as a form of curation, where the users are selecting the correct `solutions' before sharing them. Here, we want to study whether iteratively trained models can improve their planning capabilities having access only to their previously generated curated traces. We simulate the scenario just described in well-controlled environments of classical planning, and focus on the self-improving capabilities of LLMs.
First, we prompt a base
model to solve a diverse set of planning tasks, mixing together both
short-horizon and harder long-horizon planning problems. Then, we discard traces that do not solve the task, mix the remaining into the original data, and fine-tune the next-gen
model. The iteratively trained models essentially
bootstrap each other's planning capabilities: each model attempts to generate solutions to the planning tasks,
without relying on hand-designed prompts or on external planners.  The simple
plans solved in earlier generations are used as part of the training set of
later generations, allowing the model to use these simple ``building blocks'' to
solve more complicated planning problems.

In our experiments using Qwen3 4B \citep{qwen3-tr}, within five (deployment) generations, the
latest model generation achieves more than \emph{double} the base model
performance in all tested domains. Figure~\ref{fig:results-summary} summarizes
the results of our main experiment. In some cases, the performance increases by
a factor of 5. Later generations of the model can also find much longer plans
than the base model, showing that this mechanism allows for out-of-distribution
generalisation. Moreover, there is no significant difference in the average
number of reasoning tokens produced by later generations, contrasting some
results of RL fine-tuning \cite{deepseek-arxiv2025}.

We formally prove that iterative deployment is equivalent to a special case of
REINFORCE \cite{williams-jml1992} when the reward function is binary and traces
are weighted according to importance sampling.
This connection has two important implications: First, it highlights a big safety risk in the deployment of iteratively-trained models, as when the curation is done indirectly through user interactions during post-deployment, the next generation model would effectively be trained with an implicit reward function which can be difficult to control. That could have significant implications on the model's behaviour (e.g.\ the implicit reward could clash with safety training).
Second, the mechanism highlighted here can be viewed as an alternative training regime to explicit RL, preserving generalisation but using a curation mechanism instead of defining an explicit reward
function (which is often challenging for more open-ended, and less
well-specified, tasks). The curation mechanism relies on post-processing validation as a
signal: It works exclusively on the traces generated from prompts provided by
the environment after deployment (e.g., user generated interactions with the
LLM, or tool-usage traces).
Our
observations point the community to study the properties of these implicit reward functions as a
top priority, especially as such iteratively-deployed AI systems are already being used through society.

Throughout the paper we refer to the starting model simply as the base model or as the
\emph{generation 0} of our model. Analogously, we call the method deployed after
$n$ iterations of our process as the \emph{$n$-th generation} of the model.
